%% title: Maximum Likelihood Estimation
\documentclass[11pt]{article}
\usepackage{metricsstyle}

\begin{document}

\handoutheader{MLE I}{Maximum likelihood estimation: notation, invariance, consistency, asymptotic normality and the information matrix equality}{}

\runin{Context} Handed out with Lesson 14, which motivates ML with the model
$y_{t}^{\gamma}=\beta_{1}+\beta_{2}x_{2t}+u_{t}$ and introduces the Jacobian of a transformation.
Continued in \emph{MLE~II}.

\section{Introduction}

\topic{When least squares is not appropriate}
Least squares (linear or nonlinear) presumes the model can be written as
\emph{dependent variable $=$ regression function $+$ additive error}. Not every model can. Example:
\begin{equation}\tag{1}
 y_{t}^{\gamma}=\beta_{0}+\beta_{1}x_{t}+u_{t}, \qquad u_{t}\sim\text{IID}(0,\sigma^{2}).
\end{equation}
With $\gamma$ known, (1) is a regression model; with $\gamma$ to be estimated it is not, and the
sum of squares
\begin{equation}\tag{2}
 \operatorname{SSR}(\beta,\gamma)=\sum_{t=1}^{n}\bigl(y_{t}^{\gamma}-\beta_{0}-\beta_{1}x_{t}\bigr)^{2}
\end{equation}
can be driven to zero without reference to the data: if all $y_{t}>1$, take
$\gamma\to-\infty$ and $\beta_{0}=\beta_{1}=0$. Least squares cannot give sensible estimates of (1).

\topic{The idea of maximum likelihood}
Choose the parameter value $\hat\theta$ under which the observed sample has the highest
joint density, i.e.\ maximise the likelihood of the observed data.

\topic{A first example: the exponential distribution}
Let the $y_{t}$ be independent with density
\begin{equation}\tag{3}
 f(y_{t},\theta)=\theta e^{-\theta y_{t}}, \qquad y_{t}>0,\ \theta>0 .
\end{equation}
The joint density of the sample, viewed as a function of $\theta$ for the observed
$y=(y_{1},\ldots,y_{n})$, is the \emph{likelihood function}; by independence it is the product
of the individual densities:
\begin{equation}\tag{4}
 L(y,\theta)=\prod_{t=1}^{n}\theta e^{-\theta y_{t}} .
\end{equation}
Since $\log$ is strictly increasing, the maximiser of $L$ also maximises the
\emph{log-likelihood} $\ell(y,\theta)\equiv\log L(y,\theta)$, which is easier to work with:
\[ \ell(y,\theta)=\sum_{t=1}^{n}\bigl(\log\theta-\theta y_{t}\bigr)=n\log\theta-\theta\sum_{t}y_{t}. \]
First-order condition and its solution:
\[ \frac{d\ell}{d\theta}=\frac{n}{\theta}-\sum_{t}y_{t}=0
   \quad\Longrightarrow\quad \hat\theta=\frac{n}{\sum_{t}y_{t}}=\frac{1}{\bar y}, \qquad
   \frac{d^{2}\ell}{d\theta^{2}}=-\frac{n}{\theta^{2}}<0, \]
so $\hat\theta$ is a maximum.

\section{Notation}

\begin{flatlist}
 \item $y_{t}$: the $t$-th observation on a random variable; $y=(y_{1},\ldots,y_{n})'$.
 \item $\theta$: a $k$-vector of parameters that determines the distribution of $y_{t}$;
       $\theta\in\Theta$, the parameter space; $\theta_{0}$ is the true value.
 \item $L_{t}(y_{t},\theta)$: the density of $y_{t}$ given $\theta$.
 \item $L(y,\theta)\equiv L(\theta)=\prod_{t=1}^{n}L_{t}(y_{t},\theta)$: the likelihood function.
 \item $\ell_{t}(y_{t},\theta)=\log L_{t}(y_{t},\theta)$: the contribution of observation $t$ to the
       log-likelihood.
 \item $\ell(y,\theta)\equiv\ell(\theta)=\sum_{t=1}^{n}\ell_{t}(y_{t},\theta)$: the log-likelihood function.
 \item $g(y,\theta)\equiv g(\theta)$: the \emph{gradient} or \emph{score} vector ($k\times1$), with
       $g_{i}(\theta)=\partial\ell(\theta)/\partial\theta_{i}$.
 \item $G(y,\theta)\equiv G(\theta)$: the $n\times k$ \emph{CG matrix} (``contributions to the
       gradient''), with $G_{ti}(\theta)=\partial\ell_{t}(\theta)/\partial\theta_{i}$; its $t$-th row is $G_{t}$
       ($1\times k$). Summing the rows gives the score:
       \[ g(\theta)=G'(\theta)\,\iota, \qquad \iota=(1,\ldots,1)'\ (n\times1). \]
 \item $H(y,\theta)\equiv H(\theta)$: the $k\times k$ Hessian, $H_{ij}(\theta)=\partial^{2}\ell(\theta)/\partial\theta_{i}\partial\theta_{j}$.
 \item $\mathcal H^{n}(\theta)\equiv\E\bigl(\tfrac1nH(\theta)\bigr)$ and
       $\mathcal H(\theta)=\lim_{n\to\infty}\mathcal H^{n}(\theta)$.
 \item $I_{t}(\theta)$: the information in observation $t$, the $k\times k$ matrix
       \[ [I_{t}(\theta)]_{ij}=\E\bigl[G_{ti}(\theta)G_{tj}(\theta)\bigr], \qquad
          I_{t}(\theta)=\E\bigl(G_{t}'G_{t}\bigr). \]
 \item The (average) information matrix for a sample of size $n$, and its limit:
       \[ \mathcal I^{n}(\theta)=\frac1n\sum_{t=1}^{n}I_{t}(\theta)=\frac1n\,I^{n}(\theta)
          =\frac1n\,\E\bigl(G'G\bigr), \qquad
          \mathcal I(\theta)=\lim_{n\to\infty}\mathcal I^{n}(\theta). \]
\end{flatlist}
$I_{t}$ measures the expected information in the $t$-th observation and $I^{n}=n\mathcal I^{n}$
that in the whole sample. Like $I_{t}$, $\mathcal I^{n}$ and $\mathcal I$ are symmetric and, in
general, positive semidefinite. Orders: $\mathcal I^{n}=O(1)$ while $I^{n}=O(n)$ (some authors call
$I^{n}$ the information matrix).

\topic{Two definitions of the ML estimator}
\begin{flatlist}
 \item[] \textbf{Type 1.} $\hat\theta$ is an ML estimator if
         $\ell(y,\hat\theta)\ge\ell(y,\theta^{*})$ for all $\theta^{*}\in\Theta$ (a global maximiser).
 \item[] \textbf{Type 2.} $\hat\theta$ is a root of the \emph{likelihood equations}
         $g(y,\hat\theta)=0$.
\end{flatlist}

\section{Properties of maximum likelihood}

Invariance, consistency, asymptotic normality and asymptotic efficiency.

\subsection{Invariance}

Let $\eta(\cdot)$ be a set of $k$ functions mapping $\theta$ one-to-one onto a new parameter
vector $\phi=\eta(\theta)$, and write the log-likelihood in terms of $\phi$ as
$\ell'(\phi)\equiv\ell\bigl(\eta^{-1}(\phi)\bigr)$. If $\ell(\hat\theta)>\ell(\theta^{*})$ for all
$\theta^{*}\neq\hat\theta$, then with $\hat\phi=\eta(\hat\theta)$ and $\phi^{*}=\eta(\theta^{*})$,
\[ \ell'(\hat\phi)=\ell(\hat\theta)>\ell(\theta^{*})=\ell'(\phi^{*})
   \qquad\text{for all }\phi^{*}\neq\hat\phi . \]
So maximising $\ell'(\phi)$ gives $\hat\phi=\eta(\hat\theta)$: the two sets of estimates are
equivalent and attain the same maximised value. How the model is parametrised does not
affect the estimates.

\topic{Consequence: ML is generally biased}
Suppose $\E(\hat\theta)=\theta$. For nonlinear $\eta$, in general (Jensen)
\[ \E(\hat\phi)=\E\bigl(\eta(\hat\theta)\bigr)\neq\eta\bigl(\E(\hat\theta)\bigr)=\eta(\theta)=\phi, \]
so $\hat\phi$ is biased. ML estimators cannot be unbiased in every parametrisation.
The choice of parametrisation has no effect asymptotically, but matters in small samples.

\subsection{Consistency}

Under regularity conditions (the proof is difficult and omitted),
\[ \operatorname*{plim}_{n\to\infty}\hat\theta=\theta_{0}. \]

\subsection{The asymptotic distribution of the ML estimator}

\topic{Expansion of the likelihood equations}
Let $\hat\theta$ be a type 2 estimator. A first-order Taylor expansion of the
likelihood equations around $\theta_{0}$ gives
\begin{equation}\tag{5}
 0=g(\hat\theta)=g(\theta_{0})+H(\bar\theta)(\hat\theta-\theta_{0}),
\end{equation}
where $\bar\theta$ lies between $\theta_{0}$ and $\hat\theta$, chosen to make (5) exact. Solving,
\begin{equation}\tag{6}
 \hat\theta-\theta_{0}=-H^{-1}(\bar\theta)\,g(\theta_{0}),
\end{equation}
and rescaling so that every factor is $O(1)$:
\[ \sqrt n\,(\hat\theta-\theta_{0})=-\Bigl(\frac{H(\bar\theta)}{n}\Bigr)^{-1}\frac{g(\theta_{0})}{\sqrt n}. \]

\topic{The Hessian term}
The $(i,j)$ element of $H(\bar\theta)/n$ is
$\frac1n\sum_{t}\partial^{2}\ell_{t}(\theta)/\partial\theta_{i}\partial\theta_{j}\big|_{\theta=\bar\theta}$. By a law of large numbers,
and because $\bar\theta$ lies between $\hat\theta$ and $\theta_{0}$ and is therefore consistent too,
\[ \tfrac1nH(\bar\theta)\ \longrightarrow\ \mathcal H^{n}(\bar\theta)\ \longrightarrow\ \mathcal H(\theta_{0}), \]
so
\begin{equation}\tag{7}
 \sqrt n\,(\hat\theta-\theta_{0})\ \overset{a}{=}\ -\mathcal H^{-1}(\theta_{0})\,\frac{g(\theta_{0})}{\sqrt n}.
\end{equation}

\topic{The score has mean zero}
The only stochastic term on the right of (7) is $g(\theta_{0})/\sqrt n$, with typical element
$\frac1{\sqrt n}\sum_{t}\partial\ell_{t}/\partial\theta_{i}\big|_{\theta_{0}}=\frac1{\sqrt n}\sum_{t}G_{ti}(\theta_{0})$.
For each term,
\begin{align*}
 \E\Bigl[\frac{\partial\ell_{t}}{\partial\theta_{i}}\Big|_{\theta_{0}}\Bigr]
 &=\E\Bigl[\frac{1}{L_{t}(\theta_{0})}\frac{\partial L_{t}(\theta)}{\partial\theta_{i}}\Big|_{\theta_{0}}\Bigr]
  =\int\frac{1}{L_{t}(\theta_{0})}\frac{\partial L_{t}(\theta)}{\partial\theta_{i}}\Big|_{\theta_{0}}L_{t}(\theta_{0})\,dy_{t}\\
 &=\int\frac{\partial L_{t}(\theta)}{\partial\theta_{i}}\Big|_{\theta_{0}}dy_{t}
  =\frac{\partial}{\partial\theta_{i}}\int L_{t}\,dy_{t}=\frac{\partial}{\partial\theta_{i}}\,1=0. \tag{8}
\end{align*}
(8) assumes enough regularity to interchange differentiation and integration, and uses
that $L_{t}$, being a density, integrates to one. Hence $\E\bigl(G_{ti}(\theta_{0})\bigr)=0$ and
$\E\bigl(g(\theta_{0})/\sqrt n\bigr)=0$.

\topic{The variance of the score}
With mean zero, the variance of $g(\theta_{0})/\sqrt n$ is the expected outer product
\begin{equation}\tag{9}
 \frac1n\,\E\bigl(g(\theta_{0})g'(\theta_{0})\bigr),
 \qquad\text{typical element}\quad
 \frac1n\,\E\Bigl\{\Bigl(\sum_{t=1}^{n}G_{ti}(\theta_{0})\Bigr)\Bigl(\sum_{s=1}^{n}G_{sj}(\theta_{0})\Bigr)\Bigr\}.
\end{equation}
For $t\neq s$, $G_{ti}$ depends only on $y_{t}$ and $G_{sj}$ only on $y_{s}$, so by independence
\[ \E(G_{ti}G_{sj})=\E(G_{ti})\,\E(G_{sj})=0 . \] Only the $t=s$ terms survive:
\begin{equation}\tag{10}
 \frac1n\,\E\Bigl\{\Bigl(\sum_{t}G_{ti}(\theta_{0})\Bigr)\Bigl(\sum_{s}G_{sj}(\theta_{0})\Bigr)\Bigr\}
 =\frac1n\,\E\Bigl\{\sum_{t}G_{ti}(\theta_{0})G_{tj}(\theta_{0})\Bigr\},
\end{equation}
which holds only because of independence. In matrix form,
\[ \mathcal I^{n}(\theta_{0})\equiv\frac1n\,\E\bigl(G'(\theta_{0})G(\theta_{0})\bigr)
   =\frac1n\,\E\bigl(g(\theta_{0})g'(\theta_{0})\bigr). \]

\topic{The sandwich}
Define
\begin{equation}\tag{11}
 d(\theta_{0})\equiv-\mathcal H^{-1}(\theta_{0})\,\frac{g(\theta_{0})}{\sqrt n}
 \ \overset{a}{=}\ \sqrt n\,(\hat\theta-\theta_{0}).
\end{equation}
Then $\E\bigl(d(\theta_{0})\bigr)=0$, and since $\mathcal H$ is symmetric,
\begin{equation}\tag{12}
 \Var\bigl(d(\theta_{0})\bigr)
 =\bigl(-\mathcal H^{-1}(\theta_{0})\bigr)\Bigl[\tfrac1n\E(gg')\Bigr]\bigl(-\mathcal H^{-1}(\theta_{0})\bigr)
 =\mathcal H^{-1}\,\mathcal I^{n}\,\mathcal H^{-1}.
\end{equation}

\subsection{The information matrix equality}

Recall $\mathcal H^{n}_{ij}=\frac1n\E\bigl[\sum_{t}\partial^{2}\log L_{t}/\partial\theta_{i}\partial\theta_{j}\big|_{\theta_{0}}\bigr]$. From
\[ \frac{\partial\log L_{t}}{\partial\theta_{i}}=\frac{1}{L_{t}}\frac{\partial L_{t}}{\partial\theta_{i}},
   \qquad
   \frac{\partial^{2}\log L_{t}}{\partial\theta_{i}\partial\theta_{j}}
   =\Bigl[\frac{\partial^{2}L_{t}}{\partial\theta_{i}\partial\theta_{j}}L_{t}
     -\frac{\partial L_{t}}{\partial\theta_{i}}\frac{\partial L_{t}}{\partial\theta_{j}}\Bigr]\Big/L_{t}^{2}, \]
we get
\begin{equation}\tag{13}
 \frac{\partial^{2}\log L_{t}}{\partial\theta_{i}\partial\theta_{j}}
 =\frac{1}{L_{t}}\frac{\partial^{2}L_{t}}{\partial\theta_{i}\partial\theta_{j}}
  -\frac{\partial\log L_{t}}{\partial\theta_{i}}\frac{\partial\log L_{t}}{\partial\theta_{j}} .
\end{equation}
At $\theta_{0}$ the first term has expectation zero, by the argument of (8) applied twice:
\[ \E\Bigl[\frac{1}{L_{t}(\theta_{0})}\frac{\partial^{2}L_{t}}{\partial\theta_{i}\partial\theta_{j}}\Big|_{\theta_{0}}\Bigr]
   =\int\frac{\partial^{2}L_{t}}{\partial\theta_{i}\partial\theta_{j}}\Big|_{\theta_{0}}dy_{t}
   =\frac{\partial^{2}}{\partial\theta_{i}\partial\theta_{j}}\underbrace{\int L_{t}\,dy_{t}}_{=1}=0 . \]
Therefore
\[ \E\Bigl(\frac{\partial^{2}\log L_{t}}{\partial\theta_{i}\partial\theta_{j}}\Big|_{\theta_{0}}\Bigr)
   =-\E\Bigl(\frac{\partial\log L_{t}}{\partial\theta_{i}}\frac{\partial\log L_{t}}{\partial\theta_{j}}\Big|_{\theta_{0}}\Bigr)
   =-\bigl[I_{t}(\theta_{0})\bigr]_{ij}, \]
and, averaging over $t$ and taking limits,
\[ \mathcal H^{n}(\theta_{0})=-\mathcal I^{n}(\theta_{0}), \qquad \mathcal H(\theta_{0})=-\mathcal I(\theta_{0}). \]

\topic{Three equal forms of the asymptotic covariance}
Substituting into (12),
\[ \Var\bigl(\sqrt n(\hat\theta-\theta_{0})\bigr)
   \ \overset{\text{(A)}}{=}\ \mathcal H^{-1}\mathcal I\,\mathcal H^{-1}
   \ \overset{\text{(B)}}{=}\ -\mathcal H^{-1}
   \ \overset{\text{(C)}}{=}\ \mathcal I^{-1}. \]
(A) $=$ (B) $=$ (C) because the model is assumed correctly specified and everything is
evaluated at $\theta_{0}$; that is what the information matrix equality needs.

\topic{Conclusion: asymptotic normality}
$g(\theta_{0})/\sqrt n$ is a scaled sum of independent mean-zero terms, so a central limit theorem
gives
\begin{align*}
 \sqrt n\,(\hat\theta-\theta_{0}) &\ \overset{a}{\sim}\ N\bigl(0,\,-\mathcal H^{-1}(\theta_{0})\bigr), \tag{14}\\
 \sqrt n\,(\hat\theta-\theta_{0}) &\ \overset{a}{\sim}\ N\bigl(0,\,\mathcal I^{-1}(\theta_{0})\bigr), \tag{15}\\
 \sqrt n\,(\hat\theta-\theta_{0}) &\ \overset{a}{\sim}\ N\bigl(0,\,\mathcal H^{-1}(\theta_{0})\,\mathcal I(\theta_{0})\,\mathcal H^{-1}(\theta_{0})\bigr). \tag{16}
\end{align*}
Under correct specification all three coincide. Form (16), the sandwich, is built from
(12) without the information matrix equality, so it is the one to use if that equality
is in doubt.

\clearpage
\begin{transcriptnotes}
 \item Source: \emph{MLE1.pdf}, ten photographed pages of the handwritten handout headed
       ``(I) Maximum Likelihood Estimation''. Equation numbers (1)--(16) are the handout's.
       The text is restated and the derivations are laid out step by step; the section
       headings follow the handout (``1.~Introduction'', ``Notation'', ``3.~Properties''),
       with ``Notation'' numbered 2. No date. Transposes ($^{T}$ in the handout) are primes.
 \item Page 1: \textbf{Corrected:} the handout says SSR $\to0$ by $\gamma\to-\infty$ ``if all the $y_{t}$
       were greater than zero''; $y_{t}^{\gamma}\to0$ as $\gamma\to-\infty$ needs $y_{t}>1$ (for
       $0<y_{t}<1$ it explodes). The degeneracy remains for any positive $y_{t}$: at $\gamma=0$,
       $\beta_{0}=1$, $\beta_{1}=0$ the SSR is exactly zero (see Lesson 14, Section 5).
 \item Page 2: the left edge is cut off in the photo (``[s]ingle unknown parameter'',
       ``[n] observations'', ``[t]he likelihood'', ``[d]ensity''); the missing letters are
       restored from context. The second-order condition is written $d^{2}\ell/d\theta^{2}<0$ with
       its value $-n/\theta^{2}$ \textbf{added}; $\hat\theta=1/\bar y$ is \textbf{added}.
 \item Page 4: \textbf{Corrected:} the handout writes $I_{t}\equiv\E(G_{t}G_{t}^{T})$ ``since $G_{t}$ is
       $1\times k$''; for a $1\times k$ row that product is a scalar. The $k\times k$ matrix is
       $\E(G_{t}'G_{t})$, consistent with $[I_{t}]_{ij}=\E(G_{ti}G_{tj})$ and $\mathcal I^{n}=\frac1n\E(G'G)$.
       The handout's two scripts for the information matrix ($\mathcal I^{n}$ averaged,
       $I^{n}=n\mathcal I^{n}$ total) are kept as calligraphic vs.\ roman.
 \item Page 5: \textbf{Corrected:} invariance is written
       $\ell'(\hat\phi)\equiv\ell(\eta(\hat\theta))$; since $\eta$ maps $\theta$ to $\phi$, the
       reparametrised log-likelihood is $\ell'(\phi)=\ell(\eta^{-1}(\phi))$, so
       $\ell'(\hat\phi)=\ell(\hat\theta)$. The handout's $n(\theta)$ is read as $\eta(\theta)$. The
       reference to Jensen's inequality is \textbf{added}. The sentence running from page 5
       to page 6 reads ``parametrization has no asymptotic effects but matters in small samples''.
 \item Page 6: (5) is a mean-value expansion; for a vector $g$ the point $\bar\theta$ can differ
       from row to row, which does not affect the argument. ``Its $(i,j)$th element'' of
       $H(\bar\theta)/n$ is typeset as in the handout.
 \item Page 8: the sentence explaining (10) (``the product will involve taking expectations of
       terms that are independent since $L_{t}$ is independent of $L_{s}$, and hence are zero'') is
       restated precisely: cross terms vanish because they are products of independent
       mean-zero factors. The left side is written $\mathcal I^{n}(\theta)$ but is evaluated at
       $\theta_{0}$.
 \item Page 9: (12) is written $-\mathcal H^{-1}[\frac1n\E(gg')]-\mathcal H^{-1}$, the minus signs
       meant as factors; it is typeset as the product
       $(-\mathcal H^{-1})\,\mathcal I^{n}(-\mathcal H^{-1})$.
       \textbf{Corrected:} the first term of (13) was written $\frac{1}{L_{t}}\partial^{2}\ell_{t}/\partial\theta_{i}\partial\theta_{j}$;
       it is $\frac{1}{L_{t}}\partial^{2}L_{t}/\partial\theta_{i}\partial\theta_{j}$ (the density, not its log), as the
       next line of the handout uses.
 \item Page 10: the labels (A), (B), (C) are the handout's. The final remark on when to
       prefer the sandwich (16) is \textbf{added}.
 \item Original page map: p.\,1 introduction, model (1)--(2), idea of ML; p.\,2 exponential
       example (3)--(4); p.\,3 notation through $\mathcal H^{n}$; p.\,4 information matrices, type 1
       definition; p.\,5 type 2, invariance and bias; p.\,6 consistency, expansion (5)--(6),
       Hessian term; p.\,7 (7)--(8), mean of the score; p.\,8 (9)--(11), variance of the score;
       p.\,9 (12)--(13), information matrix equality; p.\,10 $\mathcal H=-\mathcal I$, (A)--(C), (14)--(16).
\end{transcriptnotes}

\end{document}
